\section{Open problems and research frontiers}
\label{sec:open}

This section states the unresolved technical problems of molten glass AM as explicit problem statements, each with the evidence from the corpus that motivates it, the physical mechanism, the current best mitigation and what a solution would have to demonstrate.

\subsection{Problem 1: Devitrification and phase purity}
\textbf{Statement.} Determine the composition-, temperature- and time-dependent crystallisation risk of glass held in the working range in crucibles, nozzles and multi-layer thermal cycles, and design process windows and feedstocks that guarantee an amorphous product.

\textbf{Evidence and mechanism.} Every silicate glass is metastable between its liquidus and \(\Tg\); nucleation and growth are fastest at intermediate viscosities (\(10^{4}\)--\(10^{7}\)\,dPa\,s), exactly where deposition happens. Crucible-fed printers hold tens of kilograms of glass at the working point for days to weeks (\SIrange{2}{3}{week} campaigns in G3DP2 \citep{inamura2018}), and every layer re-heats the layer below through the crystallisation range. Borosilicates are known to be prone to devitrification during extrusion \citep{klein2015}, and fused silica devitrifies to cristobalite above roughly \SI{1100}{\celsius} in the presence of alkali or moisture, so the laser-melted fused-silica routes operate above \SI{1900}{\celsius} for seconds and quench through the crystallisation range in milliseconds. The powder-bed papers report fully amorphous parts by XRD for soda-lime \citep{datsiou2021} and comparable XRD for DED deposits and substrate \citep{spirrett2023}; the crucible-fed papers report none. The absence of reported devitrification is therefore not evidence of its absence: soda-lime studio glasses are formulated for a long working range and low liquidus, and the risk will rise as the community moves to borosilicate, aluminosilicate and recycled compositions with contaminants.

\textbf{Required.} Time--temperature--transformation data for candidate compositions at the relevant heating rates; in-line detection (Raman or optical scattering) of early crystallisation; and nozzle designs that avoid stagnant, slowly cooled glass films on ceramic surfaces, which are the classic nucleation sites in glass forming.

\subsection{Problem 2: Optical transmittance, wavefront distortion, bubbles and layer-seam optics}
\textbf{Statement.} Quantify and then control the four optical defects of printed glass---bulk transmission loss, refractive-index inhomogeneity (striae), bubble entrapment and periodic surface relief at layer seams---to the level required for a stated application class (window, lens, waveguide).

\textbf{Evidence and mechanism.} No paper in the corpus reports a spectrophotometric transmission curve for a crucible-fed printed part; transparency is documented photographically and by caustic projection \citep{klein2015,inamura2018}. For laser-filament parts, transmittance and index homogeneity were measured for fused quartz \citep{luo2018quartz}, a transmitted-beam camera showed a homogeneous wavefront without bubbles or boundary layers in LGD specimens \citep{lgd2024homog}, and index homogeneity was compared with commercial optical glass only for the indirect DIW route \citep{destino2018}. Bubbles arise from three sources: gas dissolved in the melt released on reheating (``reboil'', observed spectroscopically above a threshold temperature \citep{luo2017heat}), gas trapped between an incoming filament and the previous layer when wetting is incomplete (the ``wiggly line'' regime \citep{kth2022bonding}), and inter-particle porosity in powder routes (\SI{12}{\percent} \citep{datsiou2021}). Striae arise from incomplete homogenisation of the melt in gravity-fed crucibles and from the thermal history difference between filament core and surface in laser routes. Layer seams act as cylindrical lenslets: a wall of \SI{10}{mm} cords is a periodic refractive surface that focuses light into caustics---attractive for art and lighting \citep{brun2017,commons_caustics} but fatal for imaging optics unless the faces are ground and polished, as was done for all mechanical and optical test samples \citep{klein2015,mit2024masonry}.

\textbf{Required.} Standardised reporting of transmission (\SIrange{400}{2000}{nm}), haze and transmitted wavefront error for printed coupons; a model linking seam relief amplitude to cord temperature, surface tension and dwell time (the visco-capillary healing time \(\eta a/\sigma\) of \cref{sec:physics}); and fining strategies (thermal soak, low-viscosity fining agents, vacuum) compatible with a printer crucible.

\subsection{Problem 3: Thermal shock, stress accumulation and polarimetric stress mapping}
\textbf{Statement.} Predict, measure and minimise the three-dimensional residual stress field of a printed glass part as a function of geometry, thermal architecture and annealing schedule, with a validated non-destructive measurement.

\textbf{Evidence and mechanism.} \Cref{eq:stress} shows that \SI{100}{\celsius} of through-thickness gradient during vitrification leaves \SI{40}{MPa} in soda-lime. Cold-printed G3DP samples had one fifth of the flexural strength of hot-printed ones and visible radial stress \citep{klein2015}; substrate cracking is the dominant failure of powder-fed DED on thin float glass \citep{spirrett2023}; and the coaxial LGD head exists because asymmetric heating of a fibre produced ``very high residual stress'' \citep{lgd2021coaxial}. Fused silica is forgiving; every other glass is not. Yet the only stress measurements in the corpus are qualitative polariscope readings \citep{klein2015} and an absence-of-birefringence statement \citep{kth2022bonding}. Photoelastic tomography and scattered-light polariscopy can deliver quantitative 3D stress in complex glass articles \citep{aben2008}, and a coupled thermo-mechanical phase-transformation framework has been proposed for glass AM, but no printed-glass study has yet combined a thermal-history model with a measured stress field.

\textbf{Required.} Thermo-viscoelastic simulation (structural relaxation with the Tool--Narayanaswamy--Moynihan model) of a multi-layer build; validation by scattered-light polariscopy on witness coupons printed alongside parts; and design rules for wall thickness and print time versus chamber temperature.

\subsection{Problem 4: Nozzle wear, chemical corrosion and seal lifetime above \SI{1200}{\celsius}}
\textbf{Statement.} Extend the continuous operating life of the glass-contact components of a molten extrusion printer beyond the two-to-three-week campaigns of alumina nozzles, and enable borosilicate and aluminosilicate processing at \SIrange{1250}{1400}{\celsius}.

\textbf{Evidence and mechanism.} Molten alkali silicates dissolve alumina and silica refractories by alkali diffusion into the glassy phase, exudation and pore formation; fused-cast AZS and high-zirconia refractories resist best but still wear and must be replaced \citep{zircorr2025}. Contact corrosion, not motion accuracy, set the campaign length of G3DP2 \citep{inamura2018}, and the crucible and refractory cost was an explicit target of the G3DP3 redesign \citep{mit2024masonry}. The glass industry's answer for fibre bushings and stirrers is platinum--rhodium (80/20 or 90/10) which operates to about \SI{1600}{\celsius} and, alloyed or dispersion-strengthened with zirconia, forms bushing nozzle tips \citep{pgm2010}; Corning's application specifies exactly this for an induction-heated printer crucible \citep{patent_corning2021}. The barriers are cost (a \SI{10}{mm} bore, \SI{50}{mm} long Pt--Rh nozzle is tens of grams of platinum-group metal), oxidation and volatilisation of rhodium above \SI{1200}{\celsius} in air, and thermal-expansion mismatch at the ceramic--metal joint. Sealing is the second problem: refractory mortar joints between crucible and nozzle, and fibreboard skirts between kiln and moving carriage, are consumables; the patent's shears and torch manifold add further wear parts.

\textbf{Required.} Life tests of candidate nozzle materials (Pt--Rh, Pt--Au dispersion alloys, iridium, dense zirconia, SiC-coated graphite under inert gas) in the specific glass; nozzle geometries that avoid stagnant films; and modular, quick-change nozzle cartridges. The laser processes avoid the problem by removing the nozzle, which is why the fused-silica market has gone to them.

\subsection{Problem 5: Closed-loop monitoring and real-time slump control}
\textbf{Statement.} Close the loop from measured cord geometry and temperature to feed rate, travel speed and heat input on the layer time scale, so that bead width, layer height and overhang stability are held to specification without operator intervention.

\textbf{Evidence and mechanism.} Crucible-fed printers currently control only temperatures and travel speed; flow drifts with head and nozzle wetting, and geometry errors accumulate at convexity changes \citep{klein2015,inamura2018}. G3DP3 corrects set points between prints by comparing printed and predicted mass \citep{stern2024cgc}---a batch, not a real-time, loop. Laser-filament processes have moved further: closed-loop melt-pool temperature from a pyrometer \citep{lgd2025tempctrl} and from a thermal camera on a deterministic real-time platform \citep{dgf2026control} are demonstrated, and the latter shows that multilayer walls are impossible without it. What is missing on all platforms is geometric feedback: a measurement of cord width and height (structured-light or laser-line profilometry through the chamber window, or thermal-image edge detection) and of slump (deviation of the top layer from the toolpath), converted into travel-speed and feed corrections. Metal DED has such controllers; glass adds the difficulties of a radiating, transparent, viscous cord in a \SI{500}{\celsius} chamber.

\textbf{Required.} A validated slump model (the sag-velocity scaling of \cref{sec:physics}) as the plant model; sensing of top-layer temperature and width at \(\ge\)\SI{10}{Hz}; and model-predictive control of travel speed and nozzle temperature with the annealing chamber as a slow disturbance.

\subsection{Problem 6: Interlayer bonding, anisotropy and mechanical qualification}
\textbf{Statement.} Establish statistically qualified mechanical properties of printed glass---including the anisotropy between along-layer and across-layer directions---sufficient for structural design codes.

\textbf{Evidence.} Klein et al.\ measured a \SI{40}{\percent} difference between orientations and \SI{60}{\percent} of material strength across layers with sample counts too small for statistics \citep{klein2015}; Inamura et al.\ reported \SIrange{41}{56}{MPa} on eight cylinders and adopted a safety factor of ten \citep{inamura2018}; the masonry study tested multiples of one design and found initial-fracture strengths from \SI{3.6}{MPa} (fully printed, thick base) to \SI{42}{MPa} and ultimate strengths of \SIrange{64}{118}{MPa}, with dimensional accuracy under \SI{1}{mm} \citep{mit2024masonry}. Powder-bed parts are an order of magnitude weaker \citep{datsiou2021}, and quartz DED reaches \SI{95}{\percent} of bulk \citep{lang2024}. Four-point bending on statistically meaningful populations, Weibull analysis, and fractography of layer interfaces are the missing elements, together with surface treatments (fire polishing, chemical strengthening) whose effect on printed glass was inconsistent \citep{inamura2018}.

\subsection{Problem 7: Feedstock chemistry and recycled glass}
\textbf{Statement.} Extend the material set beyond soda-lime studio glass to container and float cullet, borosilicate, aluminosilicate and specialty glasses while maintaining a stable process.

\textbf{Evidence.} Recycled container and float glass have been printed on G3DP3 after VFT-based viscosity matching, with \SI{295}{\celsius} working ranges nearly identical to studio glass; contamination sensitivity of post-consumer cullet was flagged as the next study \citep{stern2024cgc}. Maple Glass Printing draws its own cane from recycled cullet and must segregate colours \citep{maple_ammedia}. Borosilicate has been claimed by Micron3DP \citep{micron3dp_acers} and laser-deposited from rods \citep{froehlich2022}, but no crucible-fed borosilicate print with reported properties exists in the corpus; the required \SI{1270}{\celsius} working point (\cref{tab:vft}) exceeds current nozzle limits (Problem 4). Multi-material printing (dual crucibles, colour gradients from frit added to the melt \citep{klein2015}) and functionally graded glass are surveyed in a 2026 review \citep{multimat2026}. Molten lunar regolith extruded from a platinum--rhodium crucible under vacuum has been demonstrated as an in-situ resource utilisation route \citep{regolithreview2023,regolithvac2024}, and Evenline has stated a target of \SI{1700}{\celsius} operation for borosilicate, silica and regolith \citep{stern2024cgc}.

\subsection{Problem 8: Energy, cost and throughput}
\textbf{Statement.} Establish the unit economics of molten glass AM relative to casting, pressing and blowing at the batch sizes where geometric customisation has value.

\textbf{Evidence.} The G3DP2 thermal module consumes \SI{15}{kW} during printing at \SI{5.2}{kg/h}, i.e.\ about \SI{3}{kWh/kg} before annealing, compared with roughly \SI{1}{kWh/kg} for a modern container furnace; the cane-fed Maple machines consume a few hundred watts but at grams-per-minute rates \citep{inamura2018,maple_ammedia}. Capital costs are USD~45\,000 for a Maple~3 \citep{maple_aniwaa}, price on request for Nobula \citep{nobula_aniwaa}, and not disclosed for Evenline, which sells parts and services rather than machines \citep{evenline_web}. The masonry study argues that AM wins where every unit differs or where tooling for pressed or cast blocks cannot be amortised \citep{mit2024masonry}; no published study yet quantifies the crossover batch size. \Cref{sec:industry} returns to this.

\subsection{Problem 9: Toolpath optimisation of transparent deposited materials for target optical properties}
\label{sec:open:toolpath}
\textbf{Statement.} Given a fixed global shape and a specified optical response---a transmittance or haze map, a caustic image, a reflection or scattering signature, or a lens or waveguide function---compute the deposition toolpath of a molten-glass (or other transparent extrusion) process, i.e.\ cord layout, layer height, bead width, coiling pattern, travel-speed and temperature schedule and hence surface relief, such that the printed part realises the optical target within the geometric tolerance of the shape.

\textbf{Evidence and mechanism.} In every direct deposition route the optical response of the part is a function of the toolpath, because each cord is a cylindrical lenslet whose relief is set by bead width, layer height and the visco-capillary healing time \(\eta a/\sigma\) of the seam (\cref{sec:physics}), and because the sewing-machine instabilities of a falling viscous thread provide a catalogue of periodic surface patterns selected purely by the ratio of travel to extrusion speed \citep{brun2017}. The literature has so far recorded this dependence rather than designed with it. Klein et al.\ created caustic patterns by projecting light through G3DP prints and noted that the striations and bulges of the print texture give ``a high degree of control'' over the redirection of light \citep{klein2015,klein2015thesis}; the G3DP2 columns are described as architectural-scale lenses whose geometric variation ``can drive, limit or control light transmission, reflection and refraction'' \citep{inamura2018}, yet their geometry was optimised for load, with the caustics an emergent by-product; the accompanying \emph{Glass~II} design study stated the intention to ``reverse-engineer geometry based on designed light patterns'' as future work, for which no publication has followed \citep{inglessis2017glass2}. The two computational halves of the problem do exist in isolation. Forward simulation: Weber, Reinhart and Oxman validated Radiance photon mapping of geometrically complex printed glass against measurements \citep{weber2020photon}. Inverse simulation: Kassubeck et al.\ built a differentiable multi-spectral caustic renderer with a regulariser that uses the known deposited glass volume, in order to reconstruct the shape of laser-deposited glass from its caustic image as a process-control signal; they explicitly distinguish this ``shape from caustic'' task from the ``caustic from shape'' task of computational caustic design \citep{kassubeck2021caustics}. Caustic design itself is mature for milled or polished freeform surfaces---micro-facet and continuous-surface formulations \citep{papas2011caustics,kiser2013caustics,yue2014caustics}, optimal-transport formulations for high-contrast images \citep{schwartzburg2014caustics}, and differentiable-rendering formulations for extended light sources \citep{zhou2025caustics}---and the graphics community has likewise solved the design of projected images from perforated shells \citep{zhao2016lampshades}, light routing through printed light pipes \citep{willis2012printedoptics,pereira2014routing}, spatially varying translucency and subsurface scattering \citep{dong2010sss,papas2013translucent,elek2017scattering,brunton2018translucency,nagasawa2023translucency} and custom reflectance from fabricated microgeometry \citep{weyrich2009microgeometry}; all of these assume a fabrication process that reproduces a computed surface or voxel field, never the toolpath of a self-smoothing viscous cord. Toolpath optimisation for glass exists only with a mechanical objective: Chhadeh et al.\ showed on rod-fed soda-lime deposition that layer orientation and toolpath govern bending strength through the heat retained during deposition \citep{chhadeh2026toolpath}---the same heat-retention physics that governs seam healing and therefore optical relief. In the laser-filament routes, deposition parameters already act as optical design variables at the single-element scale: the focal length of each lens in a printed silica microlens array is set by the number of layers and the radius \citep{kth2024microlens}, deposited fibre spheres form micro-lens matrices \citep{rettschlag2020microlens}, the deposition path of a coreless fibre is the light path of a curved waveguide \citep{kranert2020waveguide}, and fibre-fed printing with temperature feedback yields spherical and aspherical germanate lenses \citep{hong2023germanate}. In the indirect direct-ink-writing route the toolpath literally encodes an optical field, the refractive-index profile of gradient-index glass, through inline ink mixing \citep{destino2018,dyllaspears2020grin}, and dedicated gradient toolpath generators with curvature-based velocity compensation have been published for it \citep{sage2026fermat}; but the variable is composition, not cord geometry, and the part is sintered and polished afterwards. Finally, for transparent polymer extrusion the toolpath--optics relation has been quantified as a scalar design-of-experiments problem: layer height, road width, raster angle, orientation, extrusion multiplier and speed against transmittance and translucency \citep{kumar2022transparency,kumar2024transparency,vochozka2024translucency,hamp2021transmittance,quero2021microfluidic}, against absorption and scattering coefficients when \emph{low} transparency is the target \citep{amendola2021diffuse}, and against birefringence \citep{ren2023anisotropy}; FDM lenses are tuned through nozzle diameter and raster angle \citep{blachowicz2021optical,heinrich2016optics,gao2023review} or made clear by coating over the layer relief \citep{rooney2024lenses}, and the direction of the printed cords has been used to engineer the reflection anisotropy of a surface \citep{anisotag2023}. No publication was found, in web, OpenAlex or Crossref searches, that closes the loop from a target optical response to the toolpath of a glass or transparent-extrusion process under a fixed shape envelope; the commercial lighting and art products of \cref{sec:industry} rely on the default layer-seam caustics.

\textbf{Required.} (i) A forward model of cord cross-section relief and its healing as a function of the toolpath variables (bead spacing, layer height, over- or under-extrusion, travel-speed modulation, coiling regime, non-planar layers), coupled to a differentiable ray or photon model of the kind already available \citep{weber2020photon,kassubeck2021caustics}. (ii) A toolpath parameterisation that spans the range from seam-free transparency (dense, fully healed cords, the glass analogue of the aligned-rectilinear, zero-cooling recipes of polymer printing) to prescribed refractive relief for caustics, diffusion or directional reflection, while respecting the molten-glass constraints of a single continuous path, nozzle-set layer height, slump and thermal history (Problems 3 and 5). (iii) An optimisation formulation---optimal transport or gradient descent on the differentiable objective---with these fabrication constraints, and a statement of the trade-off between optical target, shape tolerance and mechanical strength, since the toolpath now serves three objectives. (iv) Validation by transmitted-wavefront and caustic-image measurement on printed coupons, which would at the same time supply the geometric feedback signal that Problem 5 lacks.
